\documentclass[a4 paper,fontset=windows]{ctexart}
\usepackage{fontspec}
\usepackage{listings}
\usepackage{color}
\usepackage{graphicx}
\definecolor{dkgreen}{rgb}{0,0.5,0}
\definecolor{dkred}{rgb}{0.5,0,0}
\setmonofont{Consolas}
\usepackage{geometry}
\geometry{top=4cm,bottom=4cm,left=2.5cm,right=2.5cm}

\lstset{frame=tb,
	language=Python,
	aboveskip=3mm,
	belowskip=3mm,
	showstringspaces=false,
	columns=fixed,
	basicstyle={\small\ttfamily},
	numbers=none,
	numberstyle=\tiny\color{black},
	keywordstyle=\color{blue},
	commentstyle=\color{dkred},
	stringstyle=\color{dkgreen},
	breaklines=true,
	breakatwhitespace=true,
	tabsize=4}

	\begin{document}

	\title{\textbf{数据结构与算法K02\ \ \ \ 课堂练习}}
	\author{1900011604\ \ \ \ 张植竣\ \ \ \ 指导老师：陈斌\\北京大学物理学院天文学系}
	\date{2020年3月3日}

	\maketitle

	\section{验证\texttt{list.insert(0, item)}的时间复杂度为O(n)}
	
	\textbf{代码如下：}
	\begin{lstlisting}
	import numpy as np
	import matplotlib.pyplot as plt
	from scipy import optimize as opt
	from timeit import Timer
	
	'''
	计时实验程序1
	list.insert(0, item)：添加数据项成为列表的第一个元素
	'''

	def test1():
		return l.insert(0, -1)
	
	t1 = Timer('test1()', 'from __main__ import test1')
	x_list1 = []
	y_list1 = []
	for n in range(100000, 1100000, 100000):  # 10个数据点
		l = list(range(n))
		x_list1.append(n)
		y_list1.append(t1.timeit(number = 1000))
	for i in range(len(x_list1)):
		x_list1[i] /= 1000000
		print('%.7f'%(y_list1[i]))
	
	# 线性拟合数据点

	def f1(n, k, b):
		return k * n + b
	
	fita, fitb = opt.curve_fit(f1, x_list1, y_list1)
	print('\nk=' + str(fita[0]))
	print('b=' + str(fita[1]))
	print('t=%.2fn+%.2f\n'%(fita[0], fita[1]))
	
	# 绘图

	xmax1 = x_list1[-1]  # x轴的上界
	dx1 = x_list1[-1] / 5  # x轴的单位长度
	ymax1 = float('%.4f'%max(y_list1))  # y轴的上界
	dy1 = ymax1 / 5  # y轴的单位长度
	x1 = np.arange(0,xmax1 + dx1, 0.01) # 绘直线用
	plt.figure(figsize = (16, 10))
	plt.scatter(x_list1, y_list1, color = 'g', marker = '+')
	plt.plot(x1, f1(x1, fita[0], fita[1]), color = 'r')
	plt.title('K02-1', fontsize = 24)
	plt.xlabel("n(*10^6)", fontsize = 16)
	plt.ylabel("time(s)", fontsize = 16)
	plt.xlim((0, xmax1 + dx1))
	plt.ylim((0,ymax1+dy1))
	xaxis = np.arange(0, xmax1 + dx1, dx1)
	yaxis = np.arange(0, 2 * ymax1 + 4 * dy1, dy1)
	plt.xticks(xaxis)
	plt.yticks(yaxis)
	
	plt.show()
	\end{lstlisting}
	
	\textbf{输出结果：}
	\begin{lstlisting}
	0.0419501
	0.0979203
	0.1435118
	0.1903678
	0.2540433
	0.3134986
	0.3717452
	0.5151925
	0.5997339
	0.7364233
	
	k=0.7409057582920955
	b=-0.08105948706332189
	t=0.74n+-0.08
	\end{lstlisting}
	
	\textbf{绘制图像如下：}
	\begin{figure}[ht]
		\centering
		\includegraphics[width=0.9\linewidth]{1900011604_K02_picture1.pdf}
	\end{figure}

	\textbf{实验结果分析：}

	可以看到，运行\texttt{list.append(item)}函数$10^6$次时，数据点集中分布在直线$t=0.74n-0.08$附近，拟合出的斜率$k\approx0.74$（考虑到重复运行了$10^3$次，实际上$k\approx7.4\times10^{-4}$），说明list取值所需时间与问题的规模（即列表长度$n$）成线性关系，时间复杂度为O(n)。

	\section{验证\texttt{list.append(item)}的时间复杂度为O(1)}
	
	\textbf{代码如下：}
	\begin{lstlisting}
	import numpy as np
	import matplotlib.pyplot as plt
	from scipy import optimize as opt
	from timeit import Timer
	
	'''
	计时实验程序2
	list.append(item)：添加数据项成为列表的最后一个元素
	'''

	def test2():
	return l.append(-1)
	
	t2 = Timer('test2()', 'from __main__ import test2')
	x_list2 = []
	y_list2 = []
	for n in range(100000, 1100000, 100000):  #10个数据点
		l = list(range(n))
		x_list2.append(n)
		y_list2.append(t2.timeit(number = 100000))
	for i in range(len(x_list2)):
		x_list2[i] /= 1000000
		print('%.7f'%(y_list2[i]))
	
	# 线性拟合数据点

	def f2(n,k,b):
		return k * n + b
	
	fita, fitb = opt.curve_fit(f2, x_list2, y_list2)
	print('\nk=' + str(fita[0]))
	print('b=' + str(fita[1]))
	print('t=%.2fn+%.2f\n'%(fita[0], fita[1]))
	
	# 绘图

	xmax2 = x_list2[-1]  # x轴的上界
	dx2 = x_list2[-1] / 5  # x轴的单位长度
	ymax2 = float('%.4f'%max(y_list2))  # y轴的上界
	dy2 = ymax2 / 5  # y轴的单位长度
	x2 = np.arange(0, xmax2 + dx2, 0.01) # 绘直线用
	plt.figure(figsize = (16, 10))
	plt.scatter(x_list2, y_list2, color='g', marker='+')
	plt.plot(x2, f2(x2, fita[0], fita[1]), color='r')
	plt.title('K02-2', fontsize=24)
	plt.xlabel("n(*10^6)", fontsize=16)
	plt.ylabel("time(ms)", fontsize=16)
	plt.xlim((0, xmax2 + dx2))
	plt.ylim((0,ymax2+dy2))
	xaxis = np.arange(0, xmax2 + dx2, dx2)
	yaxis = np.arange(0, 2 * ymax2 + 4 * dy2, dy2)
	plt.xticks(xaxis)
	plt.yticks(yaxis)
	
	plt.show()
	\end{lstlisting}
	
	\textbf{输出结果：}
	\begin{lstlisting}
	0.0135574
	0.0141594
	0.0148340
	0.0146109
	0.0142911
	0.0157744
	0.0161587
	0.0134244
	0.0136601
	0.0141383
	
	k=4.919324249107837e-05
	b=0.01443381371663042
	t=0.00n+0.01
	\end{lstlisting}
	
	\textbf{绘制图像如下：}
	\begin{figure}[ht]
		\centering
		\includegraphics[width=0.9\linewidth]{1900011604_K02_picture2.pdf}
	\end{figure}
	
	\textbf{实验结果分析：}
	
	可以看到，运行\texttt{list.append(item)}函数$10^5$次时，数据点集中分布在水平线$t=0.0144s$附近，拟合出的斜率$k$几乎为0（考虑到重复运行了$10^5$次，实际上$k\approx4.92\times10^{-10}$），说明list取值所需时间与问题的规模（即列表长度$n$）无关，时间复杂度为O(1)。

	\end{document}